In this lecture, we review the linear algebra that the quantum part of the course rests on. We read
the product $Ax$ as a combination of the columns of $A$, define linear independence, the column rank
and the row rank of a matrix, and show that a nonzero vector orthogonal to a set of independent vectors
is independent of them.

\section{The Product \texorpdfstring{$Ax$}{Ax} as a Combination of Columns}

In the previous lecture we saw that $y = Ax$ is a linear combination of the columns of $A$, with the
entries of $x$ as coefficients. We work this out on the example
\[A = \begin{bmatrix} 2 & 1 & 3 \\ 3 & 1 & 4 \\ 5 & 7 & 12 \end{bmatrix} \in \mathbb{R}^{3 \times 3},
    \qquad x = \begin{bmatrix} x_1 \\ x_2 \\ x_3 \end{bmatrix},\]
whose columns we call $v_1 = (2, 3, 5)^{T}$, $v_2 = (1, 1, 7)^{T}$ and $v_3 = (3, 4, 12)^{T}$. Then
\begin{equation}\label{eq:columns-0917}
    Ax = \begin{bmatrix} 2x_1 + x_2 + 3x_3 \\ 3x_1 + x_2 + 4x_3 \\ 5x_1 + 7x_2 + 12x_3 \end{bmatrix}
    = x_1 \begin{bmatrix} 2 \\ 3 \\ 5 \end{bmatrix} + x_2 \begin{bmatrix} 1 \\ 1 \\ 7 \end{bmatrix}
    + x_3 \begin{bmatrix} 3 \\ 4 \\ 12 \end{bmatrix}
    = x_1 v_1 + x_2 v_2 + x_3 v_3.
\end{equation}
In this example the third column is the sum of the first two, $v_3 = v_1 + v_2$, because
$2 + 1 = 3$, $3 + 1 = 4$ and $5 + 7 = 12$. Substituting this into \cref{eq:columns-0917} gives
\begin{equation}\label{eq:twocolumns-0917}
    Ax = x_1 v_1 + x_2 v_2 + x_3 (v_1 + v_2) = (x_1 + x_3)\, v_1 + (x_2 + x_3)\, v_2,
\end{equation}
so every output $Ax$ is already a combination of $v_1$ and $v_2$ alone.

The matrix $A$ defines a linear map $T \colon \mathbb{R}^3 \to \mathbb{R}^3$, $T(x) = Ax$. In general, an
$m \times n$ matrix $A$ with columns $a_1, \dots, a_n$ defines the linear map
$T \colon \mathbb{R}^n \to \mathbb{R}^m$, $T(x) = Ax = x_1 a_1 + \dots + x_n a_n$, and its \emph{image}
is the set of all outputs,
\[\operatorname{Im}(A) = \set{Ax \given x \in \mathbb{R}^n} = \operatorname{span}(a_1, \dots, a_n).\]
The second equality is exactly the statement that $Ax$ is a combination of the columns: the image of
$A$ is its \emph{column space}.

\section{Linear Independence and Rank}

\begin{definition}[Linear Independence]
    Vectors $v_1, \dots, v_k$ are \textbf{linearly independent} if
    \[c_1 v_1 + c_2 v_2 + \dots + c_k v_k = 0 \quad\Longrightarrow\quad c_1 = c_2 = \dots = c_k = 0,\]
    i.e., the only combination of them that gives the zero vector is the trivial one. Otherwise they are
    \textbf{linearly dependent}. If $v_1, \dots, v_k$ are linearly independent, we say that a vector $v$
    is \emph{independent of} them if $v, v_1, \dots, v_k$ are linearly independent.
\end{definition}

\begin{definition}[Column Rank and Row Rank]
    The \textbf{column rank} $\operatorname{rank}_c(A)$ of a matrix $A$ is the maximum number of linearly
    independent columns of $A$. The \textbf{row rank} $\operatorname{rank}_r(A)$ is the maximum number of
    linearly independent rows of $A$.
\end{definition}

We use one fact about a subspace $W$ spanned by $r$ vectors $z_1, \dots, z_r$: any $r + 1$ vectors
$y_1, \dots, y_{r+1}$ of $W$ are linearly dependent. Indeed, write $y_j = \sum_{l=1}^{r} t_{lj} z_l$. The
homogeneous system $\sum_{j=1}^{r+1} t_{lj} c_j = 0$ ($1 \leq l \leq r$) has $r$ equations in $r + 1$
unknowns, so it has a solution $(c_1, \dots, c_{r+1}) \neq 0$ (see below), and then
$\sum_{j} c_j y_j = \sum_{l} \big(\sum_{j} t_{lj} c_j\big) z_l = 0$. Hence two bases of $W$ cannot have
different sizes (the larger one would consist of more vectors than span $W$, and so be dependent); the
common size is the \emph{dimension} $\dim(W)$.

That a homogeneous system of $r$ linear equations in more than $r$ unknowns has a nonzero solution
follows by induction on $r$. For $r = 0$ there are no equations and any nonzero vector is a solution.
For $r \geq 1$, if no equation contains $c_1$, set $c_1 = 1$ and all other unknowns to $0$. Otherwise
solve one equation that contains $c_1$ for $c_1$ and substitute into the others: this leaves $r - 1$
equations in the remaining unknowns, of which there are more than $r - 1$, so by induction they have a
nonzero solution, and the solved equation then gives the value of $c_1$. The column rank is the dimension of the image:

\begin{lemma}[Column Rank Is the Dimension of the Image]\label{lem:colbasis-0917}
    Let $A$ be an $m \times n$ matrix with $r = \operatorname{rank}_c(A) \geq 1$. Then any $r$ linearly
    independent columns of $A$ form a basis of $\operatorname{Im}(A)$. In particular,
    \[\operatorname{rank}_c(A) = \dim(\operatorname{Im}(A)).\]
\end{lemma}

\begin{proof}
    Let $w_1, \dots, w_r$ be independent columns of $A$ and let $a$ be any column of $A$. By the
    definition of $r$, the vectors $a, w_1, \dots, w_r$ are dependent (trivially so if $a$ is one of the
    $w_l$): $\beta a + \beta_1 w_1 + \dots + \beta_r w_r = 0$ with coefficients that are not all zero.
    Here $\beta \neq 0$, since otherwise this would be a nontrivial relation among $w_1, \dots, w_r$.
    Hence $a = -\frac{1}{\beta}(\beta_1 w_1 + \dots + \beta_r w_r) \in \operatorname{span}(w_1, \dots, w_r)$.
    So every column of $A$, and therefore every combination of columns, lies in
    $\operatorname{span}(w_1, \dots, w_r)$: the independent vectors $w_1, \dots, w_r$ span
    $\operatorname{Im}(A)$ and form a basis of it.
\end{proof}

% Not marked on the board; stated without proof. Page 1, rows 2-4: "# ind cols = rank_c(A) =
% dim(Im(A))", "# ind rows rank_r(A)", "[1 0 1], [0 1 1], [2 1 3] = 2[1 0 1] + 1[0 1 1]".
\begin{exercise}\label{ex:example-rank-0917}
    For the matrix $A$ of \cref{eq:columns-0917}, find $\operatorname{rank}_c(A)$,
    $\operatorname{Im}(A)$ and $\operatorname{rank}_r(A)$, and express every row of $A$ as a combination
    of $[1\ 0\ 1]$ and $[0\ 1\ 1]$.
\end{exercise}

\begin{wsolution}
    \emph{Columns.} The columns $v_1, v_2$ are independent: $a v_1 + b v_2 = 0$ reads $2a + b = 0$,
    $3a + b = 0$, $5a + 7b = 0$; subtracting the first equation from the second gives $a = 0$, and then
    $b = 0$. All three columns are dependent, since $v_1 + v_2 - v_3 = 0$. Hence
    $\operatorname{rank}_c(A) = 2$, and by \cref{lem:colbasis-0917},
    \[\operatorname{Im}(A) = \operatorname{span}(v_1, v_2), \qquad \dim(\operatorname{Im}(A)) = 2.\]

    \emph{Rows.} Let $u_1 = [1\ 0\ 1]$ and $u_2 = [0\ 1\ 1]$, so that $a u_1 + b u_2 = [a\ \ b\ \ a + b]$.
    Then
    \[[2\ 1\ 3] = 2 u_1 + 1 u_2, \qquad [3\ 1\ 4] = 3 u_1 + 1 u_2, \qquad [5\ 7\ 12] = 5 u_1 + 7 u_2.\]
    So all three rows lie in $\operatorname{span}(u_1, u_2)$, and any three vectors in the span of two
    vectors are dependent: $\operatorname{rank}_r(A) \leq 2$. The first
    two rows are independent: $a\,[2\ 1\ 3] + b\,[3\ 1\ 4] = 0$ gives $a + b = 0$ from the second entry
    and $2a + 3b = 0$ from the first, so $b = 0$ and $a = 0$. Hence $\operatorname{rank}_r(A) = 2$.
\end{wsolution}

In the example the two ranks agree. This is not a coincidence.

% Not marked on the board; stated without proof. Page 1, rows 2-3: both rank_c(A) and rank_r(A) are
% defined for the same matrix and computed to be equal in the example.
\begin{exercise}\label{ex:rank-equal-0917}
    Show that $\operatorname{rank}_r(A) = \operatorname{rank}_c(A)$ for every matrix $A$.
\end{exercise}

The solution is \cref{thm:rank-equal-0917}.

\begin{theorem}[Row Rank Equals Column Rank]\label{thm:rank-equal-0917}
    For every $m \times n$ matrix $A$, $\operatorname{rank}_r(A) = \operatorname{rank}_c(A)$. This common
    value is called the \textbf{rank} of $A$.
\end{theorem}

\begin{proof}
    We first show $\operatorname{rank}_r(A) \leq \operatorname{rank}_c(A)$. Let
    $r = \operatorname{rank}_c(A)$. If $r = 0$, then every column of $A$ is zero (a single nonzero vector
    is independent), so $A = 0$ and $\operatorname{rank}_r(A) = 0$ as well. Let $r \geq 1$, and let
    $w_1, \dots, w_r$ be independent columns of $A$; write $w_{il}$ for the $\nth{i}$ entry of $w_l$, so
    that $W = (w_{il})$ is the $m \times r$ matrix with columns $w_1, \dots, w_r$. By
    \cref{lem:colbasis-0917}, every column $a_j$ of $A$ is a combination of them: there are numbers
    $s_{lj}$ with
    \[a_j = \sum_{l=1}^{r} s_{lj} \, w_l \qquad (1 \leq j \leq n).\]
    Let $S = (s_{lj})$, an $r \times n$ matrix with rows $\rho_1, \dots, \rho_r$. Taking the $\nth{i}$
    entry on both sides, the entries of $A$ are $a_{ij} = \sum_{l=1}^{r} w_{il} \, s_{lj}$, that is,
    $A = WS$. Now fix $i$ and let $j$ run from $1$ to $n$: the same equation says that
    \[\text{row $i$ of $A$} \;=\; \sum_{l=1}^{r} w_{il} \, \rho_l.\]
    So every row of $A$ lies in $\operatorname{span}(\rho_1, \dots, \rho_r)$. Any $r + 1$ vectors in the
    span of $r$ vectors are dependent, so $A$ has at most $r$ independent rows:
    $\operatorname{rank}_r(A) \leq r = \operatorname{rank}_c(A)$.

    This holds for every matrix, in particular for the transpose $A^{T}$, whose rows are the columns of
    $A$ and whose columns are the rows of $A$. Hence
    \[\operatorname{rank}_c(A) = \operatorname{rank}_r(A^{T}) \leq \operatorname{rank}_c(A^{T})
        = \operatorname{rank}_r(A),\]
    and the two inequalities together give $\operatorname{rank}_r(A) = \operatorname{rank}_c(A)$.
\end{proof}

For the example, with $w_1 = v_1$ and $w_2 = v_2$, the rows of $S$ are exactly the two row vectors
$[1\ 0\ 1]$ and $[0\ 1\ 1]$ from above.

\section{Orthogonality}

The results of this section hold over $\mathbb{R}$ and over $\mathbb{C}$.
For $a, b \in \mathbb{R}^n$ the inner product is $\langle a, b \rangle = a \cdot b = \sum_{i=1}^{n} a_i b_i$.
For $a, b \in \mathbb{C}^n$ it is $\langle a, b \rangle = \sum_{i=1}^{n} \overline{a_i} \, b_i$, which is
linear in $b$. In both cases $\langle a, a \rangle = \sum_{i=1}^{n} \abs{a_i}^2 = \norm{a}^2$, which is
$0$ only for $a = 0$.

\begin{definition}[Orthogonality]
    Two vectors $a$ and $b$ are \textbf{orthogonal} if $\langle a, b \rangle = 0$; for real vectors,
    $a \cdot b = 0$.
\end{definition}

% Not marked on the board; stated without proof. Page 1, rows 5-7: "v is orthogonal to v_1, v_2, ...,
% v_k (ind vectors) then v is ind of v_1, v_2, ..., v_k", followed only by the definition of independence.
\needspace{12\baselineskip} % keep the exercise and the start of its solution on one page
\begin{exercise}\label{ex:orth-ind-0917}
    Let $v_1, \dots, v_k$ be linearly independent vectors. Show that if a nonzero vector $v$ is
    orthogonal to $v_1, v_2, \dots, v_k$, then $v$ is independent of $v_1, v_2, \dots, v_k$.
\end{exercise}

The solution is \cref{thm:orth-ind-0917}; the other reading of the line on the board is
\cref{cor:pairwise-orth-0917}.

\begin{theorem}[Orthogonal Implies Independent]\label{thm:orth-ind-0917}
    Let $v_1, \dots, v_k$ be linearly independent vectors in $\mathbb{R}^n$ or $\mathbb{C}^n$, and let
    $v \neq 0$ be orthogonal to each of them. Then $v, v_1, \dots, v_k$ are linearly independent. In
    particular, $v \notin \operatorname{span}(v_1, \dots, v_k)$.
\end{theorem}

\begin{proof}
    Suppose that $c \, v + c_1 v_1 + \dots + c_k v_k = 0$. Taking the inner product of $v$ with both
    sides, and using linearity in the second argument and $\langle v, v_i \rangle = 0$ for all $i$,
    \[0 = \langle v, \; c \, v + c_1 v_1 + \dots + c_k v_k \rangle
        = c \, \langle v, v \rangle + \sum_{i=1}^{k} c_i \, \langle v, v_i \rangle
        = c \, \norm{v}^2.\]
    Since $v \neq 0$, we have $\norm{v}^2 > 0$, so $c = 0$. The relation becomes
    $c_1 v_1 + \dots + c_k v_k = 0$, and the independence of $v_1, \dots, v_k$ gives
    $c_1 = \dots = c_k = 0$. So all coefficients vanish and $v, v_1, \dots, v_k$ are independent. If we
    had $v = \sum_{i} \alpha_i v_i$, then $1 \cdot v - \sum_{i} \alpha_i v_i = 0$ would be a nontrivial
    relation, so $v \notin \operatorname{span}(v_1, \dots, v_k)$.
\end{proof}

\paragraph{Remark.} The assumption $v \neq 0$ is needed: the zero vector is orthogonal to every vector,
but no set of vectors that contains it is independent ($1 \cdot 0 = 0$ is a nontrivial relation).

If $v_1, \dots, v_k$ are instead only assumed to be nonzero and pairwise orthogonal, then the implication
$c_1 v_1 + \dots + c_k v_k = 0 \Rightarrow c_1 = \dots = c_k = 0$ becomes the conclusion:

\begin{corollary}[Orthogonal Vectors are Independent]\label{cor:pairwise-orth-0917}
    Let $v_1, \dots, v_k$ be nonzero vectors that are pairwise orthogonal, i.e.,
    $\langle v_i, v_j \rangle = 0$ for all $i \neq j$. Then $v_1, \dots, v_k$ are linearly independent.
\end{corollary}

\begin{proof}
    Suppose that $c_1 v_1 + \dots + c_k v_k = 0$ and fix $j$. Taking the inner product of $v_j$ with
    both sides, all terms with $i \neq j$ vanish, and
    \[0 = \sum_{i=1}^{k} c_i \, \langle v_j, v_i \rangle = c_j \, \norm{v_j}^2.\]
    Since $v_j \neq 0$, this gives $c_j = 0$. This holds for every $j$, so the vectors are independent.
\end{proof}

In particular, the vectors of an orthonormal set, such as the basis states of a quantum system, are
always linearly independent, and $n$ of them in an $n$-dimensional space form a basis.
